\chapter{Conclusion}
\label{chap:conclusion}

On-board processing exists to shorten the path from acquisition to actionable information for time-critical Earth observation applications such as flood and wildfire response, where the value of a detection decays with every hour spent waiting for a ground-based product. Because a mission's own sensor data cannot exist before launch, and labeled examples of rare events accumulate slowly afterward, foundation models pre-trained on other missions' archives are an attractive way to reach useful performance with few mission-specific labels. This attractiveness is conditional, however, on a question this thesis set out to answer directly: whether such a model's learned representation actually survives contact with a sensor it was never trained on, and, if not, whether that can be fixed deliberately rather than assumed away.

The evidence gathered here says it can, but not for free. TerraMind's latent representation of real $\Phi$-sat-2 imagery does differ measurably from its representation of Sentinel-2 imagery of the same scenes, concentrated in the encoder's early layers and driven primarily by differing spectral response between the two sensors rather than by their difference in spatial resolution. Naively fine-tuning on physics-simulated degraded data does not close this gap; an explicit training signal that exploits co-registered sensor pairs directly is required, and this thesis applies that signal during the same step -- knowledge distillation -- that is separately required to compress the model onto space-qualified hardware that cannot run its original Vision Transformer architecture at all. Treating the domain-gap problem and the compute-constraint problem as one combined training objective, rather than as two sequential and independent fixes, is the thesis's central methodological position.

This work is limited to a single mission and sensor pair, and to one specific pre-trained GeoFM rather than a family of them; whether the spectral-response-driven account of the domain gap generalizes to other on-board missions, and whether TerraMind's broad pre-training is in fact more label-efficient than a mission-specific model trained from scratch, are both open questions this thesis frames precisely but does not close. The comparison against $\Phi$satNet, a bespoke GeoFM developed for the same mission by the same research group, is the immediate next step toward closing the second of these, and a natural extension of this work would repeat the analysis for other sensor pairs and other missions, or attempt to train a custom foundation model rather than adapt an existing one, to test whether the findings here are specific to TerraMind or general to the class of large, diversely pre-trained EO models it represents.
